\documentclass[10pt,letterpaper]{article}  
\usepackage[utf8]{inputenc}
\usepackage[T1]{fontenc}
\usepackage[letterpaper,margin=0.8cm]{geometry} \usepackage{xcolor}
\usepackage{hyperref}
\usepackage{titlesec}
\usepackage{enumitem}
\usepackage{parskip}

\definecolor{azul}{HTML}{2C3E6B}

\hypersetup{
colorlinks=true,
urlcolor=azul,
linkcolor=azul
}

\titleformat{\section}
{\large\bfseries\color{azul}} {}{0em}{}
[\titlerule]

\titlespacing{\section}{0pt}{6pt}{3pt}  
\setlist[itemize]{
leftmargin=*,
itemsep=0pt,  
topsep=0pt,   
parsep=0pt    
}

\setlength{\parskip}{0pt}

\pagestyle{empty}

\begin{document}

\begin{center}
{\Large\textbf{Emilio Izquierdo Montero}}  

\vspace{0.1cm}

{\normalsize Ingeniero en Sistemas Computacionales}

\vspace{0.1cm}

{\footnotesize  
Matehuala, San Luis Potosí, México \, | \,
(+52) 488 885 3863 \, | \,
\href{mailto:cucusneitor@hotmail.com}{cucusneitor@hotmail.com} \, | \,
\href{https://github.com/emili0p}{github.com/emili0p} \, | \,
\href{https://linkedin.com/in/emilio-izquierdo-91b16933a}{linkedin.com/in/emilio-izquierdo}
}
\end{center}

\vspace{-0.2cm}  

\section{Resumen Profesional}
\noindent
Ingeniero en Sistemas Computacionales con experiencia en \textbf{redes} y \textbf{Linux}. 
Líder del AWS Student Group, organizando eventos y talleres sobre computación en la nube. 
Certificado en Google Cloud y TOEFL B2. Apasionado por automatización, scripting, software libre y comunidades técnicas.

\section{Experiencia}
\noindent
\textbf{AWS Student Group Leader} \hfill Mar 2025 -- Mar 2026

\noindent
Tecnológico Nacional de México -- Campus Matehuala

\begin{itemize}
\item Organización de eventos y talleres sobre servicios AWS.
\item Coordinación de comunidad estudiantil en computación en la nube.
\item Promoción de Linux, redes y software libre en entornos académicos.
\end{itemize}

\section{Proyectos Destacados}
\noindent
\textbf{FRN} (C) -- Lenguaje de programación inspirado en Python y Lisp.

\noindent
\textbf{Taged} (Rust) -- Editor TUI estilo Vim para metadatos musicales.

\noindent
\textbf{WLB} (Bash) -- Scripts para automatización y administración de Linux.

\noindent
\textbf{KafkaExample + HadoopExample} (Python/Jupyter) -- Procesamiento de datos con Kafka y Hadoop.

\section{Educación}
\noindent
\textbf{Ingeniería en Sistemas Computacionales} \hfill 2022 -- 2026

\noindent
Tecnológico Nacional de México -- Campus Matehuala

\section{Certificaciones}
\noindent
\textbf{TOEFL ITP (B2)} \hfill Dic 2025 -- Dic 2027

\noindent
\textbf{Google Cloud Computing Foundations} \hfill Nov 2025

\section{Habilidades Técnicas}
\noindent
\textbf{Lenguajes:} Rust, Python, Bash, C/C++, Scala, TypeScript.

\noindent
\textbf{Cloud/Big Data:} AWS, Google Cloud, Kafka, Hadoop.

\noindent
\textbf{Otros:} Linux, Git, Redes, Scripting, Software Libre.

\noindent
\textbf{Idiomas:} Español (nativo), Inglés B2 (TOEFL ITP).

\end{document}
